\paraid{p-61c139e831ec}
The Gromov--Hausdorff space is a global analogue of a hyperspace.
For a fixed metric space,
the hyperspace of its nonempty compact subsets is equipped with the
Hausdorff distance.
In the Gromov--Hausdorff space
$\GHSpace$,
we consider the isometry classes of all nonempty compact metric spaces
and allow the ambient metric space to vary.
The ordinary unpointed
Gromov--Hausdorff distance
$\GHDistance$
is the infimum of the Hausdorff distances between isometric copies
in all common ambient metric spaces.



\paraid{p-f1859d6a167a}
Curtis
\cite[Theorem~1.6]{Curtis1980}
proves,
in particular,
that the hyperspace of nonempty compact subsets of a Banach space,
equipped with the Hausdorff metric of its norm,
is an absolute retract and hence an absolute neighborhood retract.
This result for a fixed ambient space motivates the analogous question
for the global space
$\GHSpace$.
Antonyan
\cite[p.~2]{Antonyan2020Euclidean}
asks whether
$\GHSpace$
is an absolute retract.
In this part,
we answer this question affirmatively by proving the following theorem.

\begin{theorem}\label{thm:part-iii-main}
\paraid{p-11b6e24f7baa}
The space
$\GHSpace$
is an absolute retract for all metrizable spaces.
\end{theorem}



\paraid{p-4df38d2e7f6c}
As a related result,
for each positive integer
$\FiniteSize$,
Antonyan identifies the subspace represented by compact subsets of
$\RealNumbers^\FiniteSize$
with the Hilbert cube minus a point
\cite[Corollary~5.3]{Antonyan2021Euclidean}.
Hyperspaces also appear in the study of Gromov--Hausdorff geodesics.
M\'emoli and Wan
\cite[arXiv v2, Theorem~1]{MemoliWan2023}
prove that every Gromov--Hausdorff geodesic can be realized as a Hausdorff
geodesic of compact subsets of one compact metric space.

\paraid{p-1cc29943ea5e}
For metrics on a fixed metrizable space,
the author proved an interpolation theorem with a uniform error bound
for prescribed metrics on a discrete family of closed subsets
\cite[Theorem~1.1]{Ishiki2026Interpolation}.
In this paper,
the domain of an extension problem maps into
$\GHSpace$,
so the underlying compact spaces vary with the parameter.
The coordinate pairs in the local models can be chosen to satisfy
the norm and coordinate convergence in \ref{item:model-continuity}
of \Cref{thm:local-models}.
We preserve their norms when placing the models in Banach spaces,
so that the partition of unity controls the error in all pairwise distances.


\paraid{p-402bb42dd02f}
To prove that
$\GHSpace$
is an ANR,
we use Hanner's domination theorem
(\Cref{thm:hanner-domination}).
Roughly speaking,
this theorem states that a separable metric space is an ANR if its identity
map can be deformed to a map factoring through an ANR by a homotopy
whose point tracks have arbitrarily small diameter.
For
$\GHSpace$,
we construct the auxiliary ANR as the quotient of the hyperspace of a
Banach space by a compact group action.
The Banach space and the group are obtained from the local models,
and the absolute retract property of this quotient follows from
Antonyan's results on equivariant hyperspaces and orbit spaces.

\paraid{p-433e89812f55}
Fix a continuous error function
$\ErrorControl\colon\GHSpace\to(0,\infty)$.
Choose local models whose errors are less than
$\ErrorControl$
on their domains,
and choose a countable locally finite partition of unity
$\{\PartitionWeight_\CoverIndex\}_{\CoverIndex\in\NonnegativeIntegers}$
subordinate to these domains.
For each
$\CoverIndex\in\NonnegativeIntegers$,
fix a local model of dimension
$\ModelDimension_\CoverIndex\geq1$
on an open set
$\ModelDomain_\CoverIndex$
containing
$\supp\PartitionWeight_\CoverIndex$.
For every
$\CoverIndex\in\NonnegativeIntegers$
and every compact metric representative
$\MetricSpace{\ComparisonCarrier}{\ComparisonMetric}$
whose isometry class belongs to
$\ModelDomain_\CoverIndex$,
choose a coordinate pair for the local model and write its point map and norm as
\[
 \FeatureCoordinates_{\CoverIndex,\ComparisonCarrier}\colon
 \ComparisonCarrier\to\RealNumbers^{\ModelDimension_\CoverIndex},
 \qquad
 \NormSymbol_{\CoverIndex,\ComparisonCarrier}\colon
 \RealNumbers^{\ModelDimension_\CoverIndex}\to[0,\infty).
\]

\paraid{p-b1d3e3825e91}
We first represent the varying local norms in fixed Banach spaces.
For each integer
$\ModelDimension\geq1$,
let
$\UnitSphere^{\ModelDimension-1}$
be the Euclidean unit sphere and set
\[
 \BanachAmbientSpace_\ModelDimension
 =\ContinuousFunctions(\UnitSphere^{\ModelDimension-1},\RealNumbers),
\]
with the uniform norm.
For every integer
$\ModelDimension\geq1$
and every norm
$\NormSymbol$
on
$\RealNumbers^\ModelDimension$,
let
$\NormSymbol\DualNorm$
denote its dual norm with respect to the Euclidean inner product.
\Cref{lem:dual-representation} constructs the linear isometric embedding
\[
 \NormEmbedding_\NormSymbol\colon
 (\RealNumbers^\ModelDimension,\NormSymbol)
 \to\BanachAmbientSpace_\ModelDimension
\]
defined,
for
$\FirstVector\in\RealNumbers^\ModelDimension$
and
$\SpherePoint\in\UnitSphere^{\ModelDimension-1}$,
by
\[
 \NormEmbedding_\NormSymbol(\FirstVector)(\SpherePoint)
 =\frac{\langle\FirstVector,\SpherePoint\rangle}
       {\NormSymbol\DualNorm(\SpherePoint)}.
\]
We take the
$\SummableNorm$
direct sum of these spaces over the local models and let the product
of the corresponding orthogonal groups act on it.
Thus we set
\[
 \BanachAmbientSpace
 =\left(\bigoplus_{\CoverIndex\in\NonnegativeIntegers}
        \BanachAmbientSpace_{\ModelDimension_\CoverIndex}\right)_{\SummableNorm},
\]
and
\[
 \ActingGroup
 =\prod_{\CoverIndex\in\NonnegativeIntegers}
   \OrthogonalGroup(\ModelDimension_\CoverIndex).
\]
Each orthogonal group acts on the corresponding space of functions
by precomposition with inverse orthogonal transformations of the sphere.
The resulting action of the compact metrizable group
$\ActingGroup$
on the separable Banach space
$\BanachAmbientSpace$
is continuous and consists of linear isometries.
Equation~\eqref{eq:dual-equivariance} makes the embeddings compatible
with orthogonal changes of the local norms and coordinates.

\paraid{p-b04298598a10}
Let
$\CompactHyperspace(\BanachAmbientSpace)$
be the space of nonempty compact subsets of
$\BanachAmbientSpace$,
with the Hausdorff metric induced by its norm,
and set
\[
 \AuxiliaryAR=\CompactHyperspace(\BanachAmbientSpace)/\ActingGroup.
\]
The absolute retract property of
$\AuxiliaryAR$
follows from two results of Antonyan.
Since
$\BanachAmbientSpace$
is connected and locally continuum-connected,
his equivariant hyperspace theorem
\cite[Proposition~3.1]{Antonyan2003West},
with the corrigendum \cite{Antonyan2006Correction},
implies that
$\CompactHyperspace(\BanachAmbientSpace)$
is an equivariant absolute retract,
or a
$\ActingGroup$-AR
(\Cref{thm:equivariant-hyperspace}).
His theorem on orbit spaces
\cite[Theorem~8 and Corollary~1]{Antonyan1990}
then implies that
$\AuxiliaryAR$
is an AR because
$\ActingGroup$
is compact with a countable basis
(\Cref{thm:orbit-retract}).
\Cref{lem:hyperspace-orbit-ar} also proves that
$\AuxiliaryAR$
is separable and metrizable.
In particular,
$\AuxiliaryAR$
is the ANR through which we will approximate the identity of
$\GHSpace$.

\paraid{p-8e3de8ba284c}
We combine the local point maps using the partition of unity.
For each compact metric representative
$\MetricSpace{\ComparisonCarrier}{\ComparisonMetric}$,
define
$\JointFeatureMap_\ComparisonCarrier\colon
\ComparisonCarrier\to\BanachAmbientSpace$
as follows.
For every
$\CoverIndex\in\NonnegativeIntegers$
and
$\ComparisonPoint\in\ComparisonCarrier$,
set
\[
 \bigl(\JointFeatureMap_\ComparisonCarrier(\ComparisonPoint)\bigr)_\CoverIndex
 =
 \begin{cases}
 \PartitionWeight_\CoverIndex\MetricSpace{\ComparisonCarrier}{\ComparisonMetric}
 \NormEmbedding_{\NormSymbol_{\CoverIndex,\ComparisonCarrier}}
   (\FeatureCoordinates_{\CoverIndex,\ComparisonCarrier}(\ComparisonPoint))
 &\text{if }\PartitionWeight_\CoverIndex
   \MetricSpace{\ComparisonCarrier}{\ComparisonMetric}>0,\\
 0&\text{otherwise}.
 \end{cases}
\]
For fixed
$\MetricSpace{\ComparisonCarrier}{\ComparisonMetric}$,
only finitely many coordinate maps are nonzero.
Thus
$\JointFeatureMap_\ComparisonCarrier$
is continuous and its image
$\CompactFeatureImage_\ComparisonCarrier
 =\JointFeatureMap_\ComparisonCarrier(\ComparisonCarrier)$
is nonempty and compact.
Changes of local bases and isometric representatives change
$\CompactFeatureImage_\ComparisonCarrier$
by elements of
$\ActingGroup$.
Thus its orbit is well defined by the input isometry class.
This construction applies to both trivial and nontrivial isometry groups,
although spaces with trivial isometry group are dense in
$\GHSpace$
by Rouyer's results
\cite[arXiv v1, Theorems~2 and~4]{Rouyer2011}.
Writing
$\OrbitProjection\colon\CompactHyperspace(\BanachAmbientSpace)\to\AuxiliaryAR$
for the orbit map,
we construct in \Cref{thm:global-domination} continuous maps
\[
 \GHSpace\xrightarrow{\ \ApproximationMap\ }
 \AuxiliaryAR
 \xrightarrow{\ \RealizationMap\ }\GHSpace.
\]
For every compact metric representative
$\MetricSpace{\ComparisonCarrier}{\ComparisonMetric}$,
the first map satisfies
\[
 \ApproximationMap\MetricSpace{\ComparisonCarrier}{\ComparisonMetric}
 =\OrbitProjection(\CompactFeatureImage_\ComparisonCarrier).
\]
The second map sends the orbit of a nonempty compact subset
$K\subset\BanachAmbientSpace$
to its isometry class with the metric induced by the norm of
$\BanachAmbientSpace$.

\paraid{p-5e91fc971616}
\Cref{thm:global-domination} also constructs a continuous homotopy
\[
 \ApproximationHomotopy\colon\GHSpace\times\UnitInterval\to\GHSpace.
\]
For every compact metric representative
$\MetricSpace{\ComparisonCarrier}{\ComparisonMetric}$,
its endpoints satisfy
\[
 \ApproximationHomotopy(\MetricSpace{\ComparisonCarrier}{\ComparisonMetric},0)
 =\MetricSpace{\ComparisonCarrier}{\ComparisonMetric},
 \qquad
 \ApproximationHomotopy(\MetricSpace{\ComparisonCarrier}{\ComparisonMetric},1)
 =\RealizationMap\ApproximationMap
   \MetricSpace{\ComparisonCarrier}{\ComparisonMetric}.
\]
For every compact metric representative
$\MetricSpace{\ComparisonCarrier}{\ComparisonMetric}$,
the point track satisfies
\begin{equation}\label{eq:part-iii-domination-3}
 \diam_{\GHDistance}
 \bigl\{\ApproximationHomotopy(\MetricSpace{\ComparisonCarrier}{\ComparisonMetric},t)
       \bigm| t\in\UnitInterval\bigr\}
 <\ErrorControl\MetricSpace{\ComparisonCarrier}{\ComparisonMetric}.
\end{equation}
For each
$\ErrorTolerance>0$,
we apply this construction with
$\ErrorControl\equiv\ErrorTolerance$.
Hanner's domination theorem
(\Cref{thm:hanner-domination})
and \eqref{eq:part-iii-domination-3} then imply that
$\GHSpace$
is an ANR for all separable metrizable spaces.
\Cref{thm:anr-category} extends this conclusion to all metrizable spaces.
Finally,
scaling the metrics contracts
$\GHSpace$
to the one-point space.
The ANR/ANE and AR/AE equivalences in \Cref{thm:retract-extensor},
together with \Cref{thm:contractible-ane},
therefore imply the absolute retract property in
\Cref{thm:part-iii-main}.

\medskip
\noindent\textbf{Dependence on the other parts.}
\paraid{p-065600b3665d}
\Cref{thm:local-models} is the input from Part~\ref{part:spectral}.
The continuity assertion \ref{item:model-continuity}
and the isometry identities \eqref{eq:input-isometry-1} and \eqref{eq:input-isometry-2}
in \Cref{thm:local-models} are used in the global construction.
The dependence on Part~\ref{part:measures} is through \Cref{thm:local-models}.
Part~\ref{part:hilbert} uses the absolute retract property established in
this part together with an independently constructed discrete
approximation to identify
$\GHSpace$
with Hilbert space.

\medskip
\noindent\textbf{Organization.}
\paraid{p-fc2cd8290ce7}
In Section~\ref{sec:prelim-topology},
we recall the equivariant hyperspace results and the domination and
extension theorems used in the proof.
In Subsection~\ref{subsec:ar-norm-representations},
we represent varying finite-dimensional norms isometrically
in a fixed Banach space,
with compatibility under orthogonal changes of coordinates.
In Subsection~\ref{subsec:ar-hyperspace-orbits},
we construct absolute retracts from hyperspaces modulo
compact group actions and compare their orbit metrics with the
Gromov--Hausdorff distance.
In Subsection~\ref{subsec:ar-global-approximation},
we use a partition of unity to combine local models into
global approximations with small homotopy tracks.
In Subsection~\ref{subsec:ar-extension},
we apply Hanner's
\Cref{thm:hanner-domination,thm:anr-category}
and use contractibility to obtain extension from
closed subsets of arbitrary metrizable spaces,
proving \Cref{thm:absolute-extensor}.

\medskip
\noindent\textbf{Conventions and notation.}
\paraid{p-97bafc9921fd}
All compact metric spaces are nonempty,
and Banach and Hilbert spaces are real.
We identify compact metric spaces with their isometry classes in
$\GHSpace$.
An isometry is a surjective isometric embedding.
For a metric space
$\MetricSpace{\BaseCarrier}{\MetricSymbol}$,
we write
$\CompactHyperspace(\BaseCarrier)$
for its nonempty compact subsets,
equipped with the Hausdorff metric
$\HausdorffDistance{\MetricSymbol}$.
For a compact space with an isometric embedding into an ambient metric space,
we identify the compact space with its image and use the restricted ambient metric.
The notation
$\IdentityMap$
denotes the identity map on its specified domain.
Sequences are indexed by
$\NonnegativeIntegers=\{0,1,2,\ldots\}$.
